\section{Introduction}
\label{sec:intro}

Translating C to Rust is a promising approach to enhancing the reliability of
legacy system programs.
%
Since C lacks language-level mechanisms to ensure memory safety, legacy software
written in C has suffered from critical security vulnerabilities caused by
memory bugs~\cite{chen2011linux, msrcblog}.
%
Rust is a modern systems programming language that addresses this issue by
providing a type system that guarantees memory safety at compile
time~\cite{matsakis2014rust, jung2017rustbelt}.
%
Thus, translating legacy code to Rust enables developers to detect previously
unknown bugs and reduce the risk of introducing new bugs~\cite{rust-curl,
li2024rust}.

Since code translation is laborious and error-prone when done manually, the
industry has developed an automatic C-to-Rust translator named
C2Rust~\cite{c2rust}.
%
It can process various real-world C codebases, producing Rust code that is
syntactically valid and semantically equivalent to the original program.
%
Software companies and open-source projects have used C2Rust to translate their
code~\cite{c2rust-huawei, zebra, qcms}.

Unfortunately, the translation produced by C2Rust is unsatisfactory because the
memory safety of the generated Rust code cannot be ensured by the Rust compiler,
which contradicts the goal of the translation.
%
The main reason is the use of C standard library (libc) functions through Rust's
foreign function interface.
%
C2Rust retains each libc function call during translation, instead of replacing
it with an equivalent function from the Rust standard library (Rust std).
%
However, foreign functions are not checked by the Rust compiler, potentially
compromising the memory safety of the entire program.

Therefore, replacing libc functions in C2Rust-generated code with proper Rust
std functions is an important problem for ensuring the safety of system programs
through translation.
%
Nevertheless, little work has been done on library replacements.
%
Most studies aimed at improving C2Rust-generated code have focused on language
features other than library functions, such as
pointers~\cite{emre2021translating, emre2023aliasing, zhang2023ownership,
wu2025genc2rust}, unions~\cite{hong2024tag}, and output
parameters~\cite{hong2024dont}.
%
One exception is the work by Hong and Ryu on the Lock
API~\cite{hong2023concrat}, but it addresses only a small subset of library
functions used in concurrent programs.

In this work, we aim to replace the I/O API in C2Rust-generated code.
%
The I/O API is one of the most important subsets of library functions, as almost
all C programs use it to interact with users (through terminals), files, or even
subprocesses, especially given C's primary role as a systems programming
language.
%
Specifically, we focus on replacing the I/O API provided by the \C{stdio.h}
header file in C with the one provided by Rust std.

The task of replacing the I/O API poses challenges due to the different designs
of the two APIs.
%
The libraries of the two languages differ not only syntactically, such as in
function names and argument order, where the conversion mapping can be
constructed manually or using existing mapping mining
techniques~\cite{zhong2010mining, nguyen2014statistical, meng2012history}.
%
They also differ semantically, primarily in two aspects: (1) types and (2) error
checking.

First, the two APIs provide different sets of types to perform I/O operations.
%
In libc, the type of a \emph{stream}---a target of each I/O operation that the
program reads from or writes to---is always \C{FILE*}.
%
In contrast, Rust std distinguishes the \emph{origins}, from which streams
originate, by defining multiple types such as \C{Stdin} and \C{File},
%
and the \emph{capabilities}, which regulate the operations that can be performed
on streams, through types such as \C{Read} and \C{Write}.
%
Therefore, replacing the I/O API requires determining the origin and capability
of each stream to assign the correct type to it.

Second, the two APIs offer different error checking mechanisms.
%
Since I/O operations may fail for various reasons,
%
% such as reaching the end of a file or running out of space on the device
%
the APIs provide ways to check whether each operation succeeded.
%
In libc, each operation silently sets an error indicator on the stream, which
the program can later inspect by calling functions like \C{ferror}.
%
This allows programs to check for errors at a different location (e.g., in
another function) from where the I/O operation was performed.
%
In contrast, each function in Rust std directly reports whether an error
occurred by returning a \C{Result} value, which is either \C{Ok} or \C{Err}.
%
As a result, replacing the I/O API requires propagating the directly returned error
indication values to the locations where the original program checks them.

\begin{figure}[t]
\begin{tikzpicture}[
  node distance=0.2cm and 0.4cm,
  auto,
  data/.style={
    rectangle,
    align=center,
    minimum height=1.8em
  },
  comp/.style={
    rectangle,
    draw,
    align=center,
    minimum height=1.8em
  },
  line/.style={draw, -Latex}
]
  \node [data, minimum height=1em] (ccode) {\scriptsize C Code};
  \node [comp, below=1.1em of ccode] (c2rust) {\scriptsize C2Rust};
  \node [data, right=of c2rust] (rustcode1) {\scriptsize Rust code\\[-0.6em]\scriptsize (libc I/O)};
  \node [comp, below=of rustcode1] (analyzer)
  {
    \scriptsize Origin and capability analysis (\cref{sec:analysis:oricap})\\[-0.6em]
    \scriptsize Error source analysis (\cref{sec:analysis:error})
  };
  \node [data, right=of analyzer] (result)
  {
    \scriptsize Streams' origins \& capabilities\\[-0.6em]
    \scriptsize Indicator-checking calls' sources
  };
  \node [comp, above=of result] (transformer) {\scriptsize Program transformation (\cref{sec:transformation})};
  \node [data, minimum height=1em, above=1.3em of transformer] (rustcode2)
    {\scriptsize Rust code\\[-0.6em]\scriptsize (Rust std I/O)};
  \path [line] (ccode) -- (c2rust);
  \path [line] (c2rust) -- (rustcode1);
  \path [line] (rustcode1) -- (transformer);
  \path [line] (transformer) -- ($(rustcode2.south)!0.1!(rustcode2.north)$);
  \path [line] ($(rustcode1.south)!0.1!(rustcode1.north)$) -- (analyzer);
  \path [line] (analyzer) -- (result);
  \path [line] ($(result.south)!0.8!(result.north)$) -- (transformer);
  \node[draw, inner xsep=3pt, inner ysep=6pt, fit=(c2rust) (rustcode1) (transformer) (analyzer) (result)] (largebox) {};
\end{tikzpicture}
  \vspace{-1em}
\caption{The workflow of the proposed approach}
  \vspace{-1.5em}
\label{fig:workflow}
\end{figure}

To address these challenges, we propose a static-analysis-based approach,
depicted in \cref{fig:workflow}.
%
Once C2Rust generates Rust code that still uses the libc I/O API, we perform two
static analyses on this code:
%
(1) \emph{origin and capability analysis}, a flow-insensitive analysis that
collects constraints on the origins and capabilities of streams and solves them,
%
and (2) \emph{error source analysis}, an interprocedural dataflow analysis that
identifies the API function calls responsible for setting indicators read by
each indicator-checking call.
%
We then transform the program to replace the I/O API, using the analysis results
to determine the type of each stream and to propagate errors.

Overall, our contributions are as follows:
%
\begin{itemize}
    %
  \item
    %
    We propose two static analyses, origin and capability analysis and error
    source analysis, to resolve the discrepancies between the I/O APIs of libc
    and Rust std (\cref{sec:analysis}).
    %
  \item
    %
    We propose program transformation techniques that replace the I/O API by
    assigning correct types to streams and propagating errors using the analysis
    results (\cref{sec:transformation}).
    %
  \item
    %
    We realize the proposed approach as a tool, Forcrat ({\bf F}\C{ILE*} {\bf
    O}peration {\bf R}eplacing {\bf C}-to-{\bf R}ust {\bf A}utomatic {\bf
    T}ranslator),
    %
    and evaluate it on 62 real-world C programs, showing that the approach is
    %
    (1) correct, with all 32 programs that have test suites passing the tests
    after transformation,
    %
    (2) efficient, analyzing and transforming 422k LOC in 14 seconds, and
    %
    (3) widely applicable, replacing 82\% of the API function calls
    (\cref{sec:evaluation}).
    %
\end{itemize}
%
We also discuss related work (\cref{sec:related}) and conclude
(\cref{sec:conclusion}).
